\cotefile{72}
\scanpage{67}

% L01 -- top left
\[
\Theta=(2-\alpha_1)\rho
\]

% L02 -- upper left and right
\deleted{\[
\begin{aligned}
\sigma\sigma'-\beta&=\underbrace{(2-\alpha_1)}\rho
& n=3 \qquad \alpha_0,\alpha_1,\alpha_2\quad\text{\unsure{paramètres}}\\
&\quad(4-\beta_1)=L_2(\beta_1)=
\end{aligned}
\]}

% L03 -- upper middle, boxed
\deleted{\[
\boxed{\sigma\sigma'-\beta=
 \underbrace{L_{n-1}(\beta_1,\ldots,\beta_{n-2})}\rho^{\deleted{\unsure{\xi}}}}
\]}

% L04 -- below box, two strokes at left
\deleted{\[
\begin{aligned}
\rule{1.1cm}{.5pt}\qquad &\overset{\scriptscriptstyle\Vert}{\zeta_1}=L_{n-1}(\,\\
\rule{1.1cm}{.4pt}\qquad &i=n-1
\end{aligned}
\]}

% L05 -- middle, above horizontal separator
\deleted{\[
\psi'\psi=\beta\,\mathrm{id}+\gamma\,\widetilde{\unsure{o|\sigma}}'\otimes\widetilde{o}
\qquad
\left\{\begin{aligned}
\beta&=\beta_0\beta_1\cdots\beta_{n-1}\\
\gamma&=L_{n-1}(\beta_1,\ldots,\beta_{n-2})
\end{aligned}\right.\quad ?
\]}

\noindent\rule{.60\linewidth}{.4pt}

% L06 -- below separator
\[
\begin{aligned}
&\deleted{\unsure{\alpha}}\\
&\Theta(a'_{n+1})=\sigma\psi'(a'_{n+1})+\gamma\\
&\langle\underbrace{\Theta(\overbrace{a'_{n+1}}^{s'_{n+1}})}_{t_{n+1}},\pi'\rangle
=\langle\Theta(\pi'),\underbrace{s'_{n+1}}_{a'_{n+1}}\rangle
=-\beta\langle\pi,s'_{n+1}\rangle=-\beta
\end{aligned}
\]

% L07 -- lower half, two columns
\[
\begin{aligned}
u&=s_0+s_1=2a_{-1}+a_0
&\deleted{f(u)=}\quad\langle u,\pi\rangle&=2\\
v&=t_n+t_{n+1}=2t_{n+1}\unsure{-}\sigma a_n
&\langle v,\pi\rangle&=-2\beta
\end{aligned}
\]

% L08 -- bottom, under one left brace
\[
\left\{\begin{aligned}
f(u)&=\deleted{f(s_0)+f}4f(\underbrace{a_{-1}}_{0})+
 \underbrace{f(a_0)}_{1}+\unsure{2}\varphi(\underbrace{a_{-1},a_0}_{-1})=-1\\
f(v)&=\underbrace{4f(t_{n+1})}_{2\beta\sigma\gamma}+
 \underbrace{\sigma^2f(a_n)}_{\beta}-
 2\sigma\varphi(\underbrace{a_n,t_{n+1}}_{\beta\sigma})
 =\beta\sigma(2\gamma-\sigma)\\[1ex]
f(s_0)&=0\\
f(t_{n+1})&=\beta\widetilde{\gamma}\widetilde{\sigma}
\end{aligned}\right.\quad\Big|\ \text{\unsure{pour mémoire}}
\]
% archive sheet number: 66
\scanpage{68}

% L01 -- top left, letterhead
% letterhead
\begin{tabular}{c}
U N I V E R S I T E\\
DES SCIENCES ET TECHNIQUES\\
DU LANGUEDOC\\
\texttt{------------}\\
Service de la Scolarité\\
MR / GJ\\
Diffusion \quad : Liste A\\
N\textsuperscript{o} : 4060
\end{tabular}

% L02 -- top right
Montpellier, le 28 Juin 1978

% L03 -- upper centre
\uline{AVIS DE SOUTENANCE D'UN DIPLOME D'ETUDES SUPERIEURES D'UNIVERSITE}

% L04 -- middle right, date and venue
Le 29 JUIN 1978 \qquad à \quad 10 Heures\\
Laboratoire de Monsieur le Professeur HALLE\\
Institut de Botanique\\
Université des Sciences et Techniques du Languedoc,

% L05 -- middle, candidate
\uline{Monsieur Ibrahima DIAITE}

% L06 -- middle, announcement
soutiendra un Diplôme d'Etudes Supérieures d'Université portant sur le sujet suivant :

% L07 -- middle, dissertation title
``Ecologie de la forêt de Basse-Casamance''.

% L08 -- lower left
\uline{Membres du jury}:\\
MM. :\\
\hspace*{1em}David W. LEE\\
\hspace*{1em}B. SAUGIER\\
\hspace*{1em}F. HALLE

% L09 -- lower right, office
Le Secrétaire Général,\\

% L10 -- lower right, signature
\handwritten{\unsure{E Siau}}\\

% L11 -- below signature
\uline{E. SIAU}
% archive sheet number: 67
\scanpage{69}

% L01 -- top, example and parameters
\[
\begin{aligned}
\text{\uline{Ex}}\qquad\underline{n=1}\qquad
\deleted{t_2=}\quad
\beta&=\beta_0=2+\unsure{\alpha_0|\alpha},\ \deleted{\unsure{4}}\rho=\deleted{\unsure{\beta_2}}4-\beta=2-\alpha\\
&\hspace{6em}\sigma=\sigma'=2\\
&\hspace{6em}\gamma=1
\end{aligned}
\]

% L02 -- upper left, brace
\[
\left\{\begin{aligned}
\beta\sigma\,\psi^{-1}&=2\beta\psi^{-1}=2\psi'+\pi\otimes\pi\\
t_2&=a_1-\beta a_{-1}\\
f(t_2)&=\beta
\end{aligned}\right.
\]

% L03 -- middle, second example
\[
\underline{n=2},\quad
\beta=\underset{\deleted{\unsure{=\alpha_0+2}}}{\beta_0}\beta_1
=\deleted{\unsure{\rho=}}(\alpha_0+2)(\alpha_1+2),\quad
\rho=-2(\alpha_0+\alpha_1),\ \sigma=2-\alpha_1,\ \sigma'=2-\alpha_0,\ \gamma=2
\]

% L04 -- lower middle
\[
\Theta=\beta\sigma\psi^{-1}=(4-\alpha_1^2)(\alpha_0+2)\psi^{-1}
=(2-\alpha_1)\psi'+2\pi\otimes\pi
\]

% L05 -- lower left
\[
t_3=2a_2+(\underbrace{2+\alpha_1^{\deleted{\unsure{1}}}}_{\beta_1})a_1+
 \underset{!}{0\cdot a_0}\ \unsure{-}\ \beta_0\beta_1a_{-1}
\]

% L06 -- foot
\[
f(t_3)=(4-\alpha_1^2)(2+\alpha_0)
\]
% archive sheet number: 68
\scanpage{70}

% L01 -- upper right, sheet upright
\mspage{\unsure{3|5}}

% L02 -- upper left, first attendance list
\uline{Etaient représentés} :

M. GOUT, par M. ROUZEYRE\\
Mme ROBIN, par Mme VERNET\\
M. THALER, par M. ASSENMACHER\\
M. KAMENKA, par M. BERTRAND\\
M. CHATELLIER, par M. FERRATON\\
M. GIRARD, par M. LABORDE\\
M. TRIAIRE, par M. RIBES\\
M. LERNER, par M. MANI

% L03 -- below proxies
\uline{Etaient absents ou excusés} :

M.M. CHARLES, CORRIU, LEFRANC, PIETRASANTA, BENOIST, SIVADE, CHICOIX, MAGOUS, Mle JUSTRABO, Mme RUSSELLO.

% L04 -- above situation heading
\uline{Assistaient à la réunion, en qualité de Directeur d'U.E.R.} :

M.M. BRUN, CANO, DURANTE, VIALLEFONT.

% L05 -- upper middle heading
\uline{SITUATION DE L'U.S.T.L.}

% L06 -- upper middle subheading
\uline{Présidence de l'Université}

% L07 -- middle, first paragraph
L'absence de candidature à la présidence de l'Université amène à un échange de vues auquel participent un certain nombre de membres du Conseil.

% L08 -- middle, torn right edge
Les statuts de l'U.S.T.L. prévoyant la nécessité d'être membre du Con\illegible{1} pour faire acte de candidature à la présidence, le Conseil admet que la démissi\illegible{1} de certains membres permettrait à de nouveaux élus d'entrer au Conseil. Parmi ce\illegible{1} nouveaux membres certaines candidatures pourraient alors se manifester.

% L09 -- middle subheading
\uline{Commissions de spécialistes}

% L10 -- above quoted motion
La motion suivante est adoptée à l'unanimité par le Conseil :

% L11 -- middle, quoted motion
``Le Conseil de l'Université des Sciences et Techniques du Languedoc exprime le désir qu'un débat soit instauré en son sein, sur la situation créée par la mise en place de nouvelles Commissions de spécialistes''.

% L12 -- below horizontal rule
\noindent\rule{\linewidth}{.4pt}

\uline{120ème séance}\\
\uline{Vendredi 3 Mars 1978}

% L13 -- lower middle heading
\uuline{CONSEIL RESTREINT AUX ENSEIGNANTS}

% L14 -- agenda
I. Approbation d'un procès-verbal\\
II. Commissions de spécialistes\\
III.Sécurité dans les laboratoires

\hfill\rule{1cm}{.4pt}\hfill\null

% L15 -- lower left, second attendance list
\uline{Etaient présents} :

M.M. CORRIU, GAUFRES, GOUT, LEFRANC, MAURIN, Mme PARIS, M.M. ROUZEYRE, THALER, Mme VERNET, M.M. DUCROT, FERRATON, RIBES.

% L16 -- lower left, proxies
\uline{Etaient représentés} :

M. AVEROUS, par M. ROUZEYRE\\
M. ASSENMACHER, par M. THALER\\
M. DUSSEAU, par M. DUCROT\\
M. PEYTAVIN, par M. RIBES

% L17 -- foot, absences
\uline{Etaient absents ou excusés} :

M.M. AMANIEU, CHARLES, PIETRASANTA, Mme ROBIN, M.M. ROUZAUD BERTRAND, Mme CASSAGNE MEJEAN, M.M. LABORDE, LERNER, SIVADE.
% archive sheet number: 69
\scanpage{71}

% L01 -- top
Supposons $\rho$ inv. \unsure{donc} \qquad $\unsure{\mathcal E|E}\simeq V\oplus\Delta_0$.

% L02 -- top, introduction of the map
On définit, pour $\lambda$ inv,

% L03 -- upper left, arrow and its value
\[
\begin{aligned}
\deleted{\unsure{s_{\lambda}}}\ p_\lambda:\lambda E&\xrightarrow{\sim}V\\
x&\longmapsto \frac1\lambda x-\underset{\displaystyle\searrow\,c\ \text{si \deleted{\unsure{cela existe}}}}{\left(\frac{\pi}{\rho}\right)}
\end{aligned}
\]

% L04 -- upper left, below map
Posons

% L05 -- upper left
\[
\widetilde{s}_i=p_1(s_i)=s_i-c=
\]

% L06 -- upper middle, brace
\[
\left\{\begin{aligned}
\widetilde{s}_0&=\deleted{\unsure{-\bigl[\frac{1}{\rho}a_0+2|a_{-1}+}}
-\frac1{\rho}\bigl[\deleted{\unsure{\beta_{n-1}}}
 \overbrace{\rho'_{n-1,1}}^{\sigma}a_0+\rho'_{n-1,2}a_1+
 \cdots+\rho'_{n-1,n-1}a_{n-2}+2a_{n-1}+a_n\bigr]\\
\widetilde{s}_i&=\widetilde{s}_0-\frac1{\rho}[a_0+\cdots+a_{i-1}]
\end{aligned}\right.
\]

% L07 -- upper middle, left qualifier
\[
\text{2 inv}\quad
\widetilde{u}=p_2(\deleted{\unsure{u}}\inserted{u})=\frac12u-c
=a_{-1}+\frac12a_0-c=\frac12(\widetilde{s}_0+\widetilde{s}_1)
=\widetilde{s}_0+\frac12a_0
\]

% L08 -- centre, open frame
\[
\boxed{-\rho\widetilde{s}_0=
 \overbrace{\rho'_{n-1,1}}^\sigma a_0+\rho'_{n-1,2}a_1+
 \cdots+\rho'_{n-1,n-1}a_{n-2}+2a_{n-1}+a_n}
\]

% L09 -- centre, second open frame and parity notes
\[
\begin{aligned}
-\rho\widetilde{u}&=-\rho\widetilde{s}_0\unsure{-}\widetilde{\rho}a_0
=(\sigma-\widetilde{\rho})a_0+\rho'_{n-1,2}a_1+\cdots+
 \rho'_{n-1,n-1}a_{n-2}+2a_{n-1}+a_n\\
&\hfill\text{\uline{$n$ pair}}\quad(\rho=2\widetilde{\rho})\\
-2\rho\widetilde{u}&=(4\widetilde{\sigma}-\rho)a_0+2\rho'_{n-1,2}a_1+\cdots+
 2\rho'_{n-1,n-1}a_{n-2}+4a_{n-1}+2a_n\\
&\hfill\text{\unsure{cas impair|car d'ailleurs}}\quad\sigma=2\widetilde{\sigma}
\end{aligned}
\]

% L10 -- lower middle, two identities at left
\[
\begin{aligned}
\widetilde{t}_i&=p_{-\beta}(t_i)=-\frac1{\beta}t_i-c\\
-\rho\widetilde{t}_i&=+\frac{\rho}{\beta}t_i+\pi
\end{aligned}
\]

% L11 -- lower middle right, three coefficients
\[
\begin{aligned}
&\text{si }\unsure{i\ge n+1|i=n+1}\\
&\text{coeff. de }a_0\text{ est }\deleted{4}-1/\rho\\
&\text{coeff. de }a_n\text{ est }-\frac{\gamma}{\beta}-\frac1\rho
=-\frac{\beta+\gamma\rho}{\beta\rho}=-\frac{\sigma\sigma'}{\beta\rho}\\
&\text{Coeff de }a_{n-1}:\ -\frac{2\gamma-\sigma}{\beta}-\frac2\rho
=\frac{-2\gamma\rho-\sigma\rho-2\beta}{\beta\rho}
=\frac{-2\sigma\sigma'-\sigma\rho}{\beta\rho}
=-\frac{\sigma}{\beta\rho}[2\sigma'+\rho]
\end{aligned}
\]

% L12 -- lower left
\[
-\rho\beta\widetilde{t}_i=\rho t_i+\beta\pi\ \deleted{\unsure{=}}
\]

% L13 -- lower part, cancelled expansion with supralinear and sublinear terms
\[
\begin{aligned}
-\rho\beta\widetilde{t}_{n+1}
&=\inserted{\rho t_{n+1}+\beta\pi}
 \deleted{\unsure{\rho t_{n+1}+\beta\pi}=(\rho\gamma+\beta)a_n+
 [\rho(2\gamma-\sigma)+2\beta]a_{n-1}+{}}\\
&\quad\deleted{\unsure{\rho\sigma a_n^{\prime *}+\beta\pi}}
 \qquad\deleted{2\sigma\sigma'-\rho\sigma=\sigma(2\sigma'-\rho)}\\
&=-\rho\sigma a_n^{\prime *}+\underbrace{\rho\gamma\pi+\beta\pi}_{\sigma\sigma'\pi}\\
&\deleted{t_{n+1}=\Theta \unsure{s'_{n+1}}}
=\sigma\bigl[\sigma'\underbrace{\pi}_{a_{n+1}^{\prime *}}-\rho a_n^{\prime *}\bigr].
\end{aligned}
\]

% L14 -- lower part, open box
\[
\boxed{-\frac{\rho\beta}{\sigma}\widetilde{t}_{n+1}
=\sigma'a_{n+1}^{\prime *}-\rho a_n^{\prime *}=\ldots}
\]

% L15 -- foot
\[
-\frac{\rho\beta}{\sigma}\widetilde{t}_i
=-\frac{\rho\beta}{\sigma}\widetilde{t}_{\unsure{n+1|n}}+
 \rho\beta\bigl(a_n+\beta_{n-1}a_{n-1}+\cdots+
 (\beta_{n-1}\cdots\beta_i)a_i\bigr)
\]
% archive sheet number: 70
\scanpage{72}

% L01 -- upper right, only title
Calculs centrés en $\unsure{o|\sigma}$\\
(si centre à dist. finie,\\
i.e. $\rho$ inv.)
% archive sheet number: 71
\scanpage{73}

% L01 -- top, bracketed and cancelled
\deleted{\uline{Choisissons} \inserted{\unsure{loc (Zar)}} section $s_0$ de $X$ admissible pour le syst. des $(\unsure{\tau_i^X|G^X|\xi_i^x})$ (déterminé mod homothétie dans $E=X\setminus C$ muni de \uline{l'origine $c$})}

% L02 -- upper third, crossed block
\deleted{(car $c\notin C$ en tt pt \inserted{i.e. $\rho$ inv} ($n$ impair et $\rho'$ inv., \deleted{\unsure{car}} ou $n$ pair \inserted{(alors $\rho=2\rho'$)} et $2$ inv)). -- \deleted{donc est} Donc on est dans les conditions suivantes (\unsure{identifions|identifier|identifiant} $E$ à $V$ grâce à l'origine $c=0$)}

% L03 -- upper third, item a
\deleted{\[
\text{a)}\quad \varphi_G^V:G\longrightarrow
\deleted{\unsure{\mathrm{Aut}(V,f)}}\mathrm O(f)
\]}

% L04 -- upper third, item b
\deleted{b) \deleted{\unsure{section $s_0$ de}} base $s_0$ de
$D_0^*=\bigcap_{1\le i\le n}H_i^{!}$
$(H_i^!=H_i\setminus H_i\cap\overline C)$}

% L05 -- left, beside item b
\margin{$a_{-1}^{\deleted{\unsure{*}}}$}

% L06 -- centre, introduction to three rows
On a \fbox{dans $V^\vee$}

% L07 -- centre, three rows
\[
\begin{aligned}
a_0^\vee&=\deleted{\unsure{\lambda_0-1}}(\lambda_0-1)a_0^*+\beta_0 a_1^*\\
a_i^\vee&=a_{i-1}^*+(\lambda_i-1)a_i^*+\beta_i a_{i+1}^*
 &1\le i\le n-1\\
a_n^\vee&=a_{n-1}^*+(\lambda_n-1)a_n^*
\end{aligned}
\]

% L08 -- left of the three rows, diagonal
% slanted
\margin{\uline{NB} $\rho$ inv. \unsure{exprime} ici que, posant $D_0=\bigcap_{1\le i\le n}H_i$, on a $a_0^\vee|D_0\ne0$ en chaque fibre (conséquence de la deuxième condition dans $2^\circ$)}

% L09 -- middle, cancelled line
\deleted{\unsure{Cela} \unsure{pour} $a_{-1}=$}

% L10 -- middle, warning in parentheses
(Attention au chgt des notations, il n'y a plus de $a_{-1},a_{-1}^*\ \ldots$)

% L11 -- below middle, starred calculation
\[
\deleted{\unsure{a'}}\quad\text{\textcircled{$*$}}\qquad
s_0\deleted{\unsure{=0-\frac1{\rho}(\zeta_0a_0+\zeta_1}}
\inserted{=\widetilde{s}_0=}
-\frac1\rho(\zeta_0a_0+\zeta_1a_1+\cdots+\zeta_n a_n)
\]

% L12 -- left, alongside coefficients
% slanted
\margin{On suppose pour simplifier maintenant les $\lambda_i=-1$ ($\tau_i$ \uline{réflexions})}

% L13 -- lower middle left
\[
\begin{aligned}
\zeta_n&=1\\
\zeta_{n-1}&=2\\
\zeta_i&=L_{n-i}(\beta_{n-1},\ldots,\beta_{i+1})
\end{aligned}
\]

% L14 -- lower middle right, square brackets
\[
\left[\begin{aligned}
\tau_i^V(s_0)&=s_0 &1\le\unsure{j|i}\le n\\
\text{i.e. }\langle s_0,a_i^\vee\rangle&=0 &1\le i\le n\\
\tau_0^V s_0&=s_0+a_0 &\text{i.e. }\langle s_0,a_0^\vee\rangle=1
\end{aligned}\right]
\]

% L15 -- below coefficients
\[
\zeta_0=L_n(\beta_{n-1},\ldots,\beta_1)=L_n(\beta_1,\ldots,\beta_{\unsure{n|n-1}})=\sigma
\]

% L16 -- bottom, item C
C) \unsure{Forme|Formes} quadratiques invariantes : \deleted{\unsure{ce sont les}} \inserted{elles forment} un \deleted{\unsure{torseur ss|torseur}} $\unsure{\mathcal O_S|G}$-Module \inserted{loc} \deleted{libre de} rg 1. On normalise par

% L17 -- bottom, left brace
\[
\left\{\begin{aligned}
f_V(a_0)&=1\quad\Longrightarrow\quad f_V(a_i)=\beta_0\beta_1\cdots\beta_{i-1}\\
\varphi(a_i,a_j)&=\deleted{\unsure{\{} }0\quad\text{si } i,j\text{ non consécutifs}\\
\varphi(a_i,a_{i+1})&=-\beta_0\cdots\beta_i
\end{aligned}\right.
\]
% archive sheet number: 72
\scanpage{74}

% L01 -- top
Reformulation des données lorsque $c\notin C$ en \unsure{tte fibre|ts fibres} i.e. $\rho$ inv ($n$ impair, $\rho'$ inv, ou $n$ pair, $\rho'$ inv. et 2 inv.)

% L02 -- upper third, A
\[
\text{A)}\quad \varphi_G:G\longrightarrow\mathrm{Aut}(X)\quad
\text{ou encore}\quad(\tau_i^X)_{0\le i\le n}
\]

% L03 -- upper third
est \unsure{remplacé} par

% L04 -- upper third, second action
\[
\varphi_G:G\longrightarrow\mathrm{Aut}(\unsure{W|V})\quad
\text{ou encore}\quad(\tau_i^V)_{0\le i\le n}
\]

% L05 -- above numbered list
Conditions

% L06 -- first condition
$1^\circ)$ les $\tau_i^V$ pseudo-réflexions (\unsure{d'où|définit|c.à.d.} $L_i\subset V_{\unsure{i}},\ H_i\subset V_i$)

% L07 -- left of second condition, diagonal
% slanted
\margin{\deleted{\unsure{déf. des fais}}}

% L08 -- second condition
\[
2^\circ)\quad V=\bigoplus_{0\le i\le n}L_i,\quad V^\vee\simeq
\bigoplus\underbrace{H_i^0}_{L'_0}\quad\text{i.e.}\quad
\bigcap_{0\le\unsure{j|i}\le n}H_i=0\quad\text{chaque fibre}
\]

% L09 -- third condition
\[
3^\circ)\quad L_i\subset H_{i+1}\quad\text{chaque fibre}\qquad0\le i\le n-1
\]

% L10 -- cancelled fourth condition
\deleted{\[
4^\circ)\quad\text{\unsure{Posant}}\quad D_0=\bigcap_{1\le i\le n}H_i,\quad\text{\unsure{il est|est}}
\]}

% L11 -- middle and lower part, B
B) La donnée de $s_0\in\unsure{E\subset X|\Gamma(X)}$ \deleted{\unsure{en}} \unsure{donnant lieu avec|donnant lieu à} $\varphi_G$ à une \unsure{réal. géom} non dég, est \deleted{\unsure{rempl}} exprimée par

% L12 -- middle, base
\[
\text{\uline{$\widetilde{s}_0$ base de}}\quad D_0=\bigcap_{1\le i\le n}H_i
\]

% L13 -- lower third
qui épingle la situation. \deleted{On obtient} [Il revient au $=$ de se donner une base \deleted{C) Les formes quadratiques} $a'_0$ de $L'_0$ (car $D_0\otimes L'_0\xrightarrow{\sim}\mathcal O_S$, par

% L14 -- lower third, boxed
\[
\boxed{\langle\widetilde{s}_0,a'_0\rangle=1}
\]

% L15 -- foot
ce qui fixe bien $a_0$ de $L_0$ par $\tau_0=\mathrm{id}_V+a'_0\otimes a_0$, et de proche en proche on détermine les $a'_i,a_i$ par $\langle a_i,a'_{i+1}\rangle=1,\qquad\tau_i=\mathrm{id}+a'_i\otimes a_i$ ]
% archive sheet number: 73
\scanpage{75}

% L01 -- top
Posant

% L02 -- top
\[
\omega_V=a_0\wedge\cdots\wedge a_n
\]

% L03 -- top
on a

% L04 -- upper third, discriminant
\[
\delta(f_V)=\underbrace{\beta_0^n\beta_1^{n-1}\cdots\beta_{n-1}}_{B}
 \rho\,\omega_V^{\otimes-2}=B\rho\,\omega_V^{\otimes-2}
\]

% L05 -- below discriminant
($f$ lisse sss $B$ inv. i.e. les $\beta_i$ inv \quad $0\le i\le1$.)

% L06 -- upper middle
Considérons alors

% L07 -- upper middle, parity cases
\[
\lambda=f_V(\widetilde{s}_0)=\deleted{\unsure{-}}+
\frac{\sigma'\rho'}{\rho^2}
=\left\{\begin{aligned}
&\left(\begin{gathered}n\text{ pair \unsure{avec|donc} }\rho=2\rho'\\(2\text{ inv!})\end{gathered}\right)
&&+\frac{\sigma'}{4\rho'}\\
&n\text{ impair }\ \rho=\rho'&&+\frac{\sigma'}{\rho'}
\end{aligned}\right.
\]

% L08 -- middle, quadric
L'\deleted{\unsure{a}} \deleted{\unsure{hyp}}hyperquadrique affine circonscrite au polyèdre régulier est alors donnée par

% L09 -- middle
\[
f_V-\lambda=0
\]

% L10 -- middle, smoothness
Pour qu'elle soit \deleted{affine} \inserted{\unsure{projectivement|respectivement}} lisse il f. et s. que les $\beta_i$ et $\sigma'$ soient inv (c'est donc le cas \uline{\unsure{birégulier|singulier}})

% L11 -- lower third, D
(D) \uline{NB} Supposons donnés $V$, $f_V$ et $\varphi_G^V$, représentation \deleted{\unsure{ortho.|orth.}}

% L12 -- lower third, action
\[
G\longrightarrow \mathrm O(f_{\unsure{V|W}})
\]

% L13 -- below action
$1^\circ$ transformant les $\tau_i$ en réflexions \uline{orth.} et satisfaisant de plus $2^\circ)$ $3^\circ)$ plus haut.

% L14 -- lower third, revised wording
Ceci \uline{implique} \deleted{\unsure{satisfaisant de plus}},

% L15 -- above foot, replacement condition
\[
\deleted{\unsure{4^\circ)}}\quad
\bigl[\deleted{f\mid(D_0=\bigcap_{1\le i\le n}H_i)}
\inserted{f\mid L_0}\ne0\quad\text{en \unsure{tte|ts} fibre.}
\]

% L16 -- lower left, diagonal, double-headed pointer to condition
% slanted
\margin{$\longleftrightarrow$ (i.e. $\tau_0$ est une réflexion \unsure{str.} orth. de $V$) \deleted{\unsure{Ceci} Cette condition est-elle une \unsure{conséq.} des autres ?}}

% L17 -- foot, continuation
Alors \unsure{ça|$s_0$} provient \unsure{loc|doc|donc} fppf d'une situation comme plus haut. Il suffit de choisir une section \inserted{$a_0$} de $D_0$ telle que $f_V(a_0)=1$
% archive sheet number: 74
\scanpage{76}

% L01 -- top, parenthesis continued
(L'indétermination est donc section d'un rev. principal de $S$ de \unsure{gpe|groupe} $\mu_2$, défini mod él. de $\mu_2(S)$ \ldots). Donc $s_0$ est alors donné par la formule $(*)$.

% L02 -- upper third
(Rien n'empêche d'ailleurs que $f_V(s_0)=0$)

% L03 -- upper middle, struck heading
\deleted{(E) Supposons $n$ impair}

% L04 -- upper middle
Si on se donne un hom

% L05 -- upper middle, homomorphism
\[
\varphi_G^C:G\unsure{-|\longrightarrow}\mathrm{Aut}(C,Q_C)
\]

% L06 -- middle, square brackets
[$C$ fibré proj. de rg $n$, $Q_C$ \inserted{hyper}quadrique,]

% L07 -- above list
satisfaisant les conditions

% L08 -- first condition
1) Les $\tau_i^X$ réflexions \deleted{\unsure{aut|orth.}} \qquad ($0\le\unsure{j|i}\le n$) -- d'où $\ell_i\in\Gamma C,\ H_i$ hyperplans)

% L09 -- second condition
2) les $\ell_i$ engendrent $C$, \deleted{\unsure{les|i.e.}} $\bigcap_{0\le i\le n}H_i^{\unsure{C}}=\varnothing$

% L10 -- third condition
3) \deleted{\unsure{$\ell_i$}} $\ell_i\notin H_{i+1}$ \qquad $0\le i\le n-1$

% L11 -- fourth condition, crossed out
\deleted{\unsure{et} 4) \deleted{\unsure{un}} $\ell_0\notin Q_C$ en chaque fibre}.

% L12 -- lower left, diagonal qualification
% slanted
\margin{Pour $n\ge2$ (si $n=1$, une réflexion de la droite \unsure{pr.} $C=\mathrm P(V)$ ne se remonte pas de \uline{façon univoque} en une réfl. de $V$ \ldots]}

% L13 -- lower part
alors loc. (\unsure{ét|et}) ça provient de $(V,f_V,(\tau_i^V))$ satisfaisant les conditions de (D), $V$ déterm. \unsure{près de|mod} $\otimes L$ ($L$ inv), $f_V$ défini mod unité. On \unsure{Utilisant|utilise} \unsure{la remarque|la remarque de|alors} ; $L_0\subset V$, on \unsure{les|l'} ``\unsure{épingler|épingle|épingles}'' i.e. éliminer l'indét. $\otimes L$, en prenant $L\simeq\mathcal O_S$ i.e. section $a_0$ de $L_0$.
% archive sheet number: 75
\scanpage{77}

% L01 -- top
Et on normalise alors par $f(a_0)=1$. On trouve :

% L02 -- upper half, equivalence
La catégorie des réal. géom. de $\unsure{\Pi|T_1}$ non dég. avec centre à distance finie, est équivalente à celle des triples $(C,Q_C,\varphi_G^C)$ d'un fibré projectif de rg $n$, d'une hyperquadrique $Q_C$ de $C$, d'une rep. de $G$ dans $(C,Q_{\unsure{C|\ell}})$, satisfaisant les conditions $1^\circ)$ à $3^\circ)$.

% L03 -- left opposite third written line
\margin{\uline{NB} $n\ge2$}

% L04 -- middle, parenthesis
(\uline{NB} Si on ne se donne pas $Q_C$, \unsure{énoncé|conique} analogue, mais \deleted{\unsure{avec}} $Q_C$ est déduite \unsure{canoniquemt|canoniquement} comme unique hyperquadrique inv.
% archive sheet number: 76
\scanpage{78}

% L01 -- top edge
\[
\widetilde{s}_0=s_0-\unsure{o|\sigma|c}
\]

% L02 -- upper left, crossed diagonally
\deleted{\[
\begin{aligned}
s'_0&=\rho\widetilde{s}_0,\\
f_V(s'_0)&
\end{aligned}
\]}

% L03 -- top, evaluation
\[
f_V(\widetilde{s}_0)=\deleted{\unsure{f_{X}}}f(s_0-o)=
 \underbrace{f(s_0)}_{\substack{\Vert\\0}}+\underbrace{f(o)}_{-(\sigma'\rho')/\rho^2}-
 \underbrace{\varphi(s_0,o)}_{-\sigma/\rho}
=-\frac{\sigma'\rho'}{\rho^2}+\frac{\sigma}{\rho}
\]

% L04 -- upper third, two parity rows
\[
\left\{\begin{aligned}
n\text{ pair}\quad&\rho=2\rho',\ \sigma^{\deleted{\unsure{2}}}=\sigma'\\
n\text{ impair}\quad&\rho=\rho',\ \sigma=2\sigma'
\end{aligned}\right.\quad
\left\{\begin{aligned}
f_V(\widetilde{s}_0)&=-\tfrac14\tfrac{\sigma'}{\rho'}+\tfrac12\tfrac{\sigma'}{\rho'}=\tfrac14\tfrac{\sigma'}{\rho'}\\
f_V(\widetilde{s}_0)&=-\tfrac{\sigma'}{\rho'}+2\tfrac{\sigma'}{\rho'}=\tfrac{\sigma'}{\rho'}
\end{aligned}\right.
\]

% L05 -- upper middle
\[
\begin{aligned}
\text{Soit}\quad s'_0&=-\rho\widetilde{s}_0=\deleted{\unsure{-}}
(\underbrace{\zeta_0}_{\substack{\Vert\\\sigma}}a_0+\zeta_1a_1+\cdots+\underbrace{\zeta_n}_{\substack{\Vert\\1}}a_n)\\
f_V(s'_0)&=\rho\sigma-\rho'\sigma'=\rho'\sigma'
\end{aligned}
\]

% L06 -- upper middle, box at left
\[
\boxed{f_V(s'_0)=\rho'\sigma'}
\]

% L07 -- middle, right of box
\[
\varphi_V(\widetilde{s}_0,a_i)=\varphi_{\deleted{\unsure{V}}}(s_0-o,a_i)=
\underbrace{\varphi(s_0,a_i)}_{\substack{0\text{ si }i\ne0\\-1\text{ si }i=0}}
-\underbrace{\varphi(o,a_i)}_{0\ \forall i}
=\begin{cases}-1&i=0\\0&1\le\unsure{j|i}\le n\end{cases}
\]

% L08 -- middle left, two pairs
\[
\left\{\begin{aligned}\widetilde{s}_i&=s_i-o\\s'_i&=-\rho\widetilde{s}_i\end{aligned}\right.
\qquad\left\{\begin{aligned}
\widetilde{s}_i&=\tau_{i-1}^V\widetilde{s}_{i-1}\\
s'_i&=\tau_{i-1}^V s'_{i-1}
\end{aligned}\right.\qquad(1\le i\le n)
\]

% L09 -- middle right, bracket
\[
\begin{aligned}
\varphi_V(s'_0,a_i)&=\begin{cases}\rho&\unsure{j|i}=0\\0&1\le i\le n\end{cases}\\
\langle\widetilde{s}_0,a_i^\vee\rangle&=\langle\widetilde{s}_0,a'_i\rangle=\\
\langle s_0-o,a'_i\rangle&=\langle s_0,a'_i\rangle\\
&=\begin{cases}1&i=0\\0&1\le i\le n\end{cases}
\end{aligned}
\]

% L10 -- middle left, sums
\[
\left\{\begin{aligned}
\widetilde{s}_i&=(s_i-s_0)+(s_0-o)=\widetilde{s}_0+(a_0+\cdots+a_{i-1})\\
s'_i&=s'_0\unsure{-}\rho(a_0+\cdots+a_{i-1})
\end{aligned}\right.
\]

% L11 -- lower middle, first edge
\[
\begin{aligned}
s'_1&=s'_0-\rho a_0\\
\deleted{\unsure{f}}s'_0+s'_1&=2s'_0-\rho a_0\\
f_V(s'_0+s'_1)&=4\underbrace{f_V(s'_0)}_{\rho'\sigma'}+\rho^2
\underbrace{f_V(a_0)}_1\unsure{-}2\rho\underbrace{\varphi_V(s'_0,a_0)}_\rho
=4\rho'\sigma'-\rho^2=
\end{aligned}
\]

% L12 -- below first edge, parity split
\[
f_V(s'_0+s'_1)=\begin{cases}
n\text{ pair}:&4\rho'(\sigma'-\rho')\\
n\text{ impair}:&\rho'(4\sigma'-\rho')
\end{cases}
\]

% L13 -- lower third
Donc Pour que l'homothétie de rapport $\lambda$ transforme la sphère $\Sigma$ circonscrite aux $s'_i$ en celle circonscrite aux \unsure{transformés} de $s'_0+s'_1$, il f. et s. que

% L14 -- lower third, final ratios
\[
\lambda^2\rho'\sigma'=\deleted{\rho'(4\sigma'-\rho')}
\left\{\begin{aligned}
&4\rho'(\sigma'-\rho')\quad\text{i.e.}\quad
\lambda=2\sqrt{1-\rho'/\sigma'}\quad n\text{ pair}\\
&\deleted{\text{i.e. }\lambda=}\\
&\rho'(4\sigma'-\rho')\qquad\lambda=2\sqrt{1-\rho'/4\sigma'}\quad n\text{ imp}
\end{aligned}\right.
\]
% archive sheet number: 77
\scanpage{79}

% L01 -- top right, only title
Sous-espaces invariants
% archive sheet number: 78
\scanpage{80}

% L01 -- top, hypothesis
On suppose qu'on est sur \uline{corps $k$}.

% L02 -- upper third, theorem
\uline{Théorème} Soit $Z\subset X$ sous-espace projectif. \deleted{\unsure{Soit}} Pour que $Z$ soit invariant par $G$ i.e. par les $\tau_i$ \inserted{$(0\le i\le n)$}, il f et s. qu'on soit dans l'un des cas \unsure{suivant|suivants} :

% L03 -- upper third, a
a) $Z$ est égal à l'un des ss-espaces $\varnothing,X,\{c\},C$.

% L04 -- upper third, b
b) $\exists\ 0\le i\le n-1$ tel que $\beta_i\ (=2+\alpha_i)=0$, et on a

% L05 -- upper middle, two alternatives
\[
Z=\underbrace{H_0\cap H_1\cap\cdots\cap H_i}_{Z'_i\ (\text{de dim }n-i)}
\quad\text{ou}\quad
Z=\deleted{\unsure{L}}\underbrace{\mathrm{Lin}(h_{i+1},\ldots,h_n)}_
 {\substack{Z_i^{\deleted{\unsure{\prime|*}}}\\(\dim n-i-1)}}
\]

% L06 -- left beside alternatives
% slanted
\margin{NB $\beta_i=0\ \Longleftrightarrow\ \underbrace{Z_i\subset Z'_i}_{\text{codim }1}$}

% L07 -- left of b, isolated
\margin{\unsure{b}}

% L08 -- centre, convention
\[
\begin{aligned}
(\text{\uline{NB} On posera}\quad Z_n&=\underbrace{H_0\cap\cdots\cap H_n}_{\{c\}}
 &\text{et}\quad Z'_n&=\varnothing\ \deleted{\unsure{H_n}}.\\
 Z_{-1}&=X &Z'_{-1}&=C
\end{aligned}
\]

% L09 -- above two chains
On a (sans hyp. sur \unsure{aucun|aucune} des $\beta_i$)

% L10 -- middle, two horizontal chains with vertical inclusions
\diagram{\resizebox{\linewidth}{!}{\begin{tikzcd}[ampersand replacement=\&,column sep=2pt,row sep=1pt]
 Z_n^{\deleted{\prime}}=\varnothing \& \subset \& Z_{n-1}^{\deleted{\prime}}=\{h_n\} \& \subset \& Z_{n-2}^{\deleted{\prime}}=\unsure{\mathrm L|\mathrm{Lin}}(h_n,h_{n-1}) \& \cdots \& \subset \& Z_0^{\deleted{\prime}}=\mathrm L(h_n,\ldots,h_1) \& \subset \& Z_{-1}^{\deleted{\prime}}=C\\
 \cap \& {} \& \cap \& {} \& \cap \& {} \& {} \& \cap \& {} \& \cap\\
 Z'_n=\{c\} \& \subset \& Z'_{n-1}=H_0\cap\cdots\cap H_{n-1} \& \subset \& Z'_{n-2}=H_0\cap\cdots\cap H_{n-2} \& \cdots \& \subset \& Z'_0=H_0 \& \subset \& Z'_{-1}=X
\end{tikzcd}}}

% L11 -- below chains
Toutes les inclusions écrites sont strictes, de codim 1.

% L12 -- lower third, dimensions
\deleted{Tous} les $2(n+2)$ espaces $Z_i,Z'_i$ $( -\unsure{0|1}\le i\le n)$ sont\inserted{-ils} distincts ? \deleted{\unsure{on}} Pour une dim donnée $-1\le d\le n+1$, il y a \deleted{\unsure{au}} \inserted{un} \deleted{\unsure{seul}} seul des $Z_i,Z'_i$ de cette dim si $d=-1$ ou $d=n+1$ (savoir $Z_n^{\deleted{\prime}}$ et $Z'_{-1}$), et \inserted{exactement} deux si $0\le d\le n$, savoir $Z'_{n-d}=H_0\cap\cdots\cap H_{n-d}$ \deleted{\unsure{de}} et $Z_{n-d-1}^{\deleted{\prime}}=\mathrm{Lin}(h_n,\ldots,h_{n-d})$, \deleted{codim} et ils \deleted{\unsure{ne}} sont égaux sss $h_{n-d},h_{n-d+1}\in Z'_{n-d}$, i.e. $h_{n-d}\in H_{n-d}\cap H_{n-d-1}$ $h_{n-d+1}\in H_{n-d}$

% L13 -- lower left, framed note
\margin{\uline{NB} Si $c\notin C$ i.e. $\rho\ne0$, alors les $Z_i,Z'_i$ tous distincts, car $Z'_i\not\subset Z_j^{\deleted{\prime}}$ pour tts $i,j$}
% archive sheet number: 79
\scanpage{81}

% L01 -- top, continuation
i.e. sss on a simultanément (posant $i=n-d$)

% L02 -- top, brace and qualifications
\[
\left\{\begin{aligned}
\mu_i\ (=\lambda_i-1)&=0&&
\begin{gathered}
\text{(i.e. }2\cdot1_k=0\text{ dans le cas des réflexions\unsure{)|,}}\\
\text{i.e. } k\text{ de car. }2)
\end{gathered}\\
\beta_i=\beta_{i-1}&=0&&
\begin{gathered}
\text{(se réduit à la condition }\beta_i=0\\
\text{si }i=0,\text{ et à la condition}\\
\beta_{i-1}=0\text{ si }i=n)
\end{gathered}
\end{aligned}\right.
\]

% L03 -- upper left, abandoned headings
\deleted{\uline{Cor} \unsure{D'autre} \uline{Dém.}}

% L04 -- upper middle, corollary a
\uline{Cor.} \textcircled{a} Il y a au plus $2(n+2)$ sous-espaces projectifs invariants par $G$.

% L05 -- upper middle, b
\textcircled{b} Pour qu'il y en ait $2(n+2)$, il f et s. que $\beta_i=0$ $(0\le i\le n-1)$ et $\mu_i\ne0$ \deleted{\unsure{$\forall$}} $(0\le i\le n)$ (\unsure{cette condition|cette autre} signifie aussi, dans le cas des réflexions, que car. $k\ne2$).

% L06 -- middle, c
\textcircled{c} \deleted{Pour} Il y a au moins les quatre ss-espaces invariants $\varnothing,X,\{c\},C$. Pour que ce soient les seuls, il \deleted{\unsure{f et s.|faut et}} \inserted{suffit} que $\beta_i\ne0\qquad 0\le i\le n-1$ (ce qui implique que les $2(n+\unsure{1|2})$ espaces $Z_i,Z'_i$ sont distincts). Il faut \inserted{et s.} que $\beta_i\ne0\quad 0\unsure{\le|<}i<n-1$, et que $(\beta_0\deleted{\unsure{\ne\beta_n|=\beta_n}}=0\Rightarrow\mu_0=0)$, $(\beta_{n-1}=0\Rightarrow\mu_n=0)$ [donc, dans le cas des réflexions, que

% L07 -- lower middle, cases
\[
\left\{\begin{aligned}
&\text{si car. }k\ne2:\ \text{les }\beta_i\ne0\ \inserted{0\le i\le n-1}\quad
(\text{i.e. }\beta=\prod_i\beta_i\ne0\\
&\text{si car }k=2\qquad\text{les }\beta_i\ne0\text{ si }0<i<n-1.
\end{aligned}\right.
\]

% L08 -- lower part, struck d
\deleted{d) Pour que \deleted{\unsure{la sous-représentation|sous-espace}} \inserted{la représentation projective} de $G$ \uline{dans $C$} soit irréductible, il f et s. que $\rho\ne0$ et les $\beta_i\ne\unsure{2|0}$, i.e. que}

% L09 -- lower left, enclosed and linked to d
\margin{Si $c\notin C$ i.e $\rho\ne0$}

% L10 -- foot, new d, continues
d) les sous-espaces projectifs stables de $C$ \deleted{\unsure{et de $\{c\}$ (dans le cas $c\in C$)}} $\ne C$ \deleted{et}, de $\varnothing$ sont en corr. 1-1 avec les $0\le i\le n-1$ tels que $\beta_i=0$, \deleted{et} un tel $i$ correspondant $Z_i^{\deleted{\unsure{\prime|*}}}$ (de dim $n-1-i$). Ces sous
% archive sheet number: 80
